% CONTEXT (from the published note v1.2; statements only)
Model: SU(2) Wilson lattice gauge theory on periodic tori Lambda(n_1..n_4), weight exp(-beta sum_p (1 - 1/2 Tr U_p)), Haar measure normalised. Z(vec n) its partition function, F=log Z.
Delta(n) = 16 F(n,n,n,n) - F(2n,2n,2n,2n).
(W1) F is invariant under permutations of the four periods.
(W2) For every spatial torus sigma=(n1,n2,n3) (sides >= 3 in what follows) there is a self-adjoint, POSITIVE, trace-class operator K_sigma on the space H_inv(sigma) of gauge-invariant L^2 functions of the spatial links with Z(sigma,N)=Tr K_sigma^N for all N>=2. lambda_0>=lambda_1>=... eigenvalues with multiplicity, r_i=lambda_i/lambda_0 in [0,1], s_sigma(N)=sum_{i>=1} r_i^N, h(N)=log(1+s(N)).
Lemma (squaring). For r_i in [0,1], p(N)=2log(1+s(N))-log(1+s(2N)): (a) 0<=log(1+s(N))<=p(N), so r_i^N<=s(N)<=e^{p(N)}-1; (b) p(2N)<=2(e^{p(N)}-1)^2.
Proposition (bracket). Under (W1)-(W2), with sigma_j=((2n)^j,n^{3-j}), j=0..3 and s_j(n)=s_{sigma_j}(n): sum_j 2^{3-j} log(1+s_j(n)) <= Delta(n) <= 2 sum_j 2^{3-j} log(1+s_j(n)).
Theorem R (published). Assume (W1)-(W2) for all periods >=3. Let n>=3, Delta(n)<=2^-9, q=256 Delta(n)<=1/2, m_j=2^j n. (a) For every torus t with min t_i>=m_j: h_t(m_j) <= 2 Delta(m_j) V(t)/(31 m_j^3), V(t)=t1 t2 t3, and 256 Delta(m_j) <= q^{2^j}. (c) For every integer N>=n every excited ratio r of the cube (N,N,N) satisfies r^n<=q. (d) For every torus with min t_i>=n and max t_i<=24 min t_i, every excited ratio satisfies r^n<=theta=e^{18/31}-1. [The proof of (d) shows s_t(m)<=2^{1-p}theta<=theta^p with p=2^j, m=m_j<=a<2m.]
Example (Z_2 gauge theory at beta=infinity, d=4): F(n)=cV+3 log 2 on every rectangular torus, Delta_4(n)=45 log 2 for all n; K is a positive multiple of the rank-8 projector onto the holonomy sectors H^1(T^3;Z_2).
